\documentclass[activite]{thedocument}
\usepackage{tkz-euclide}
\startdatakeys
\source{}
\cycle{}
\competencesDuSocle{}
\connaissancesRequises{}
\competencesTravaillees{}
\enddatakeys
\begin{document}
    \begin{enumerate}
        \item Placer la fraction $\dfrac{2}{3}$ sur la demi-droite graduée
        suivante :\vskip20pt\noindent
        \begin{tikzpicture}[xscale=4]
            \coordinate[label={[label distance=10pt]-90:0}](O)at(0,0);
            \coordinate[label={[label distance=10pt]-90:1}](A)at(1,0);
            \coordinate[label={[label distance=10pt]-90:2}](B)at(2,0);
            \coordinate(C)at(1/3,0);
            \coordinate(D)at(2/3,0);
            \coordinate(E)at(4/3,0);
            \coordinate(F)at(5/3,0);
            \foreach\X/\size in{O/0.2,A/0.2,B/0.2,C/0.1,D/0.1,E/0.1,F/0.1}{
                \path(\X)--+(0,\size)coordinate(\X1);
                \path(\X)--+(0,-\size)coordinate(\X2);
                \draw(\X1)--(\X2);}
            \tkzDrawLine[add=0 and 0.3](O,B)
        \end{tikzpicture}
        \item Placer la fraction $\dfrac{8}{7}$ sur la demi droite-graduée
        suivante :\vskip20pt\noindent
        \begin{tikzpicture}[xscale=4]
            \coordinate[label={[label distance=10pt]-90:0}](O)at(0,0);
            \coordinate[label={[label distance=10pt]-90:1}](A)at(1,0);
            \coordinate[label={[label distance=10pt]-90:2}](B)at(2,0);
            \foreach\ii in{1,2,3,4,5,6,8,9,10,11,12,13}{
                \coordinate(C\ii)at(\ii/7,0);
                \path(C\ii)--+(0,0.1)coordinate(C\ii1);
                \path(C\ii)--+(0,-0.1)coordinate(C\ii2);
                \draw(C\ii1)--(C\ii2);
            }
            \foreach\X in{O,A,B}{
                \path(\X)--+(0,0.2)coordinate(\X1);
                \path(\X)--+(0,-0.2)coordinate(\X2);
                \draw(\X1)--(\X2);}
            \tkzDrawLine[add=0 and 0.3](O,B)
        \end{tikzpicture}
        \item Placer la fraction $\dfrac{14}{6}$ sur la demi droite-graduée
        suivante :\vskip20pt\noindent
        \begin{tikzpicture}[xscale=4]
            \coordinate[label={[label distance=10pt]-90:0}](O)at(0,0);
            \coordinate[label={[label distance=10pt]-90:1}](A)at(1,0);
            \coordinate[label={[label distance=10pt]-90:2}](B)at(2,0);
            \foreach\ii in{1,2,3,4,5,7,8,9,10,11,13,14,15,16}{
                \coordinate(C\ii)at(\ii/6,0);
                \path(C\ii)--+(0,0.1)coordinate(C\ii1);
                \path(C\ii)--+(0,-0.1)coordinate(C\ii2);
                \draw(C\ii1)--(C\ii2);
            }
            \foreach\X in{O,A,B}{
                \path(\X)--+(0,0.2)coordinate(\X1);
                \path(\X)--+(0,-0.2)coordinate(\X2);
                \draw(\X1)--(\X2);}
            \tkzDrawLine[add=0 and 0.3](O,B)
        \end{tikzpicture}
        \item En utilisant une graduation adaptée, placer la fraction $\dfrac{1}{2}$ sur la demi droite
        suivante : \vskip20pt\noindent
        \begin{tikzpicture}[xscale=6]
            \coordinate[label={[label distance=10pt]-90:0}](O)at(0,0);
            \coordinate[label={[label distance=10pt]-90:1}](A)at(1,0);
            \foreach\X in{O,A}{
                \path(\X)--+(0,0.2)coordinate(\X1);
                \path(\X)--+(0,-0.2)coordinate(\X2);
                \draw(\X1)--(\X2);}
            \tkzDrawLine[add=0 and 0.3](O,A)
        \end{tikzpicture}
        \item En utilisant une graduation adaptée, placer la fraction $\dfrac{2}{5}$ sur la demi droite
        suivante : \vskip20pt\noindent
        \begin{tikzpicture}[xscale=6]
            \coordinate[label={[label distance=10pt]-90:0}](O)at(0,0);
            \coordinate[label={[label distance=10pt]-90:1}](A)at(1,0);
            \foreach\X in{O,A}{
                \path(\X)--+(0,0.2)coordinate(\X1);
                \path(\X)--+(0,-0.2)coordinate(\X2);
                \draw(\X1)--(\X2);}
            \tkzDrawLine[add=0 and 0.3](O,A)
        \end{tikzpicture}
        
        \item En utilisant une graduation adaptée, placer la fraction $\dfrac{9}{8}$ sur la demi droite
        suivante : \vskip20pt\noindent
        \begin{tikzpicture}[xscale=6]
            \coordinate[label={[label distance=10pt]-90:0}](O)at(0,0);
            \coordinate[label={[label distance=10pt]-90:1}](A)at(1,0);
            \foreach\X in{O,A}{
                \path(\X)--+(0,0.2)coordinate(\X1);
                \path(\X)--+(0,-0.2)coordinate(\X2);
                \draw(\X1)--(\X2);}
            \tkzDrawLine[add=0 and 0.3](O,A)
        \end{tikzpicture}
    \end{enumerate}
\end{document}